% !TeX spellcheck = ru_RU
\section{Эксперимент}

\subsection{Условия эксперимента}

Эксперимент проводился на следующих данных:
\begin{itemize}
    \item 20 страниц русскоязычной вики «Властелин колец» (10 персонажей, 10 локаций);
    \item 15 страниц русскоязычной вики «Леди Баг» (10 персонажей, 5 локаций).
\end{itemize}
Оборудование: ноутбук с ОС Windows 11, процессор Intel Core i5, 16 ГБ ОЗУ. Python 3.13, библиотеки: requests, beautifulsoup4, pandas, spacy, natasha (экспериментально), networkx, matplotlib. База данных PostgreSQL, импорт через DataGrip 2026.

\subsection{Исследовательские вопросы}

В работе проверялись следующие вопросы:
\begin{description}
    \item[RQ1] – Возможно ли собрать тексты с фандомной вики, несмотря на защиту от ботов, и корректно извлечь сущности с помощью стандартной NER-модели?
    \item[RQ2] – Обеспечивает ли предложенная архитектура лёгкую замену фандома без изменения кода анализа?
    \item[RQ3] – Корректно ли работают SQL-запросы и визуализация графа на загруженных данных?
\end{description}

\subsection{Результаты}

\textbf{RQ1:} С использованием расширенных заголовков и пауз (3 секунды) удалось успешно загрузить все страницы LOTR и Ladybug. Ошибок 403 не возникло. Эксперимент с NER-моделями показал, что стандартная spaCy \texttt{ru\_core\_news\_sm} часто классифицирует имена персонажей как \texttt{LOC} (локации). Альтернативная библиотека Natasha была протестирована (файл \texttt{extract\_entities1.py}), но не показала значительного улучшения. В финальной версии оставлена spaCy с фильтром \texttt{['PER', 'LOC']}, что позволило выделить персонажей с типом \texttt{PER}. Например, в статье «Гэндальф» найдено 45 сущностей \texttt{PER} и 30 \texttt{LOC}. Качество распознавания оценено визуально – большинство имён распознаны правильно, но встречаются ошибки (например, краткие формы или нестандартные имена).

\textbf{RQ2:} Для проверки заменяемости был создан второй парсер \texttt{LadybugParser}. Переключение между \texttt{LotrParser} и \texttt{LadybugParser} в файле \texttt{run\_parser.py} не потребовало изменений в модуле извлечения сущностей, SQL-запросах и визуализации. Архитектура признана расширяемой.

\textbf{RQ3:} Все SQL-запросы выполнены в DataGrip и вернули осмысленные результаты. Построен граф совместных упоминаний персонажей для вселенной LOTR. После фильтрации (степень >= 5, вес ребра >= 5) граф содержит 23 вершины и 107 рёбер (исходный граф имел 256 вершин и 11896 рёбер). Крупные узлы (Фродо, Гэндальф, Саурон, Арагорн) имеют наибольшую степень, что соответствует сюжетной логике.

\subsection{Обсуждение результатов}

Полученные данные подтверждают, что система работает. Главная сложность сейчас в том, что одна и та же сущность (например, персонаж) может встречаться в тексте в разных падежах («Гэндальф», «Гэндальфа», «Гэндальфу»), и программа воспринимает их как разных персонажей. В будущем это можно исправить, научив программу приводить все формы слова к одному виду. Ещё одна проблема: модель иногда ошибочно принимает за персонажа или место то, что им не является (например, слово «Средиземье» может быть помечено как персонаж). Это исправляется дообучением модели на наших конкретных текстах. Кроме того, в перспективе можно автоматически собирать ссылки на все страницы фандомной вики, а не прописывать их вручную, как сейчас. Это сделает систему полностью автоматической и избавит от ручного ввода.
\subsection{Угрозы корректности}

\begin{itemize}
    \item \textbf{Малая выборка} – 20 и 15 страниц недостаточно для статистической значимости, но достаточно для демонстрации конвейера.
    \item \textbf{Зависимость от версий библиотек} – использована актуальная на момент разработки версия spaCy; при обновлении поведение NER может измениться.
    \item \textbf{Субъективность оценки качества NER} – проверялась вручную, без формальных метрик.
\end{itemize}

\subsection{Воспроизводимость}

Все исходные коды и SQL-скрипты находятся в репозитории \url{https://github.com/s-vera666/fandom_wiki_parser}. Для воспроизведения достаточно установить зависимости из \texttt{requirements.txt} и выполнить последовательно \texttt{run\_parser.py}, \texttt{extract\_entities.py}, затем импортировать CSV в PostgreSQL и запустить \texttt{graph.py}.